% !TEX root = ../main.tex
\section{What Instigates Compression in DNNs}\label{sec:compression}
A rebuttal paper released shortly after \cite{shwartzziv2017} written by Saxe
et al. \cite{saxe2018} provide evidence that several claims from the former
paper are wrong. Here we will focus our attentions to what has a direct causal
link to the observed compression phenomena as seen in the information plane, and
measured by the mutual information $I(X; T)$.

In \cite{shwartzziv2017} they claim that the stochasticity introduced by
SGD induces a random diffusion process which can be described by the
Fokker-Planck equation. Their argument is that this diffusion, subject to the
error constraint, maximizes the conditional entropy $H(X\mid T_i)$ i.e.
minimizes the mutual information $I(X;T_i)=H(X)-H(X\mid T_i)$, since $H(X)$ is
fixed by the data. Thus attributing the compression to the optimization method
SGD. In \cite{saxe2018} they show that compression is seen
as a result of:
\begin{itemize}
	\item \textbf{Double-saturating nonlinearities} cause the hidden activations
	to saturate as training progresses. The binned activity then collapses into
	the two extreme bins, and the measured $I(X;T)$ drops. Their tanh network
	compresses, their ReLU network does not. The same holds when the mutual
	information is estimated with kernel density and $k$-nearest-neighbour
	estimators instead of binning.
	\item \textbf{Growth in weight magnitude} make the units enter this
	saturated regime in the first place. In their minimal model of a single
	unit $h=f(w_1X)$ with $X\sim\mathcal{N}(0,1)$, $I(T;X)=H(T)$, since $T$ is
	deterministic given $X$, and they show $H(T)$
	first rises and then falls as a function of the weight $w_1$ for tanh,
	whilst it only rises for ReLU.
	\item \textbf{Genuinely task-irrelevant input dimensions} do compress, but
	this happens at the same time as the fitting and not in a separate phase
	afterwards. The information on the input as a whole can still increase.
\end{itemize}

% Your own saturation/binning sketch goes here. Uncomment once figures/ has it.
% \begin{figure}[t]
%   \centering
%   \includegraphics[width=\linewidth]{saturation}
%   \caption{}
%   \label{fig:saturation}
% \end{figure}


Subsequently they find that compression is \textit{not} a result of:
\begin{itemize}
	\item \textbf{The drift-diffusion transition} in the gradient SNR. They
	find this transition in all cases, also in networks that never compress,
	so the two phases are not causally related to the compression.
	\item \textbf{The maximum entropy argument} of \cite{shwartzziv2017}, which
	they argue does not hold theoretically. The stationary distribution it
	describes is a distribution over weights across different training runs,
	whilst $H(X\mid T)$ is about the inputs, and this gives no reason that the
	weights found in one particular run maximize $H(X\mid T)$. They do not show
	the claim to be false, but leave the burden of proof with
	\cite{shwartzziv2017}.
	\item \textbf{The stochasticity of SGD}, as they train the same networks
	with full batch gradient descent, which has no randomness in its updates,
	and still see comparable compression in the tanh networks.
\end{itemize}
where their account thus reduces to a single mechanism; the weights grow during
training, and the double-saturating nonlinearity turns this growth into a drop
in the measured $I(X;T)$. However the mechanism relies on the weights growing in
the first place, and what makes them grow is never asked, but simply taken as
given. The loss function is held fixed in all of the comparisons they make. The
permutation which would settle this, tanh trained with MSE, is never run, and
the only setting where they do use MSE is the linear network, in which the
nonlinearity is changed as well.
